% ============================================================
% BLINDED MANUSCRIPT — Indian Journal of Finance
% Times New Roman 12pt, double spaced, 1-inch margins.
% NO author name or affiliation appears in this file (double-blind review).
% ============================================================
\documentclass[12pt]{article}

\usepackage[a4paper,margin=1in]{geometry}
\usepackage{newtxtext,newtxmath}   % Times New Roman-metric text and math
\usepackage{setspace}
\usepackage{array}
\usepackage{booktabs}
\usepackage{ragged2e}
\usepackage{titlesec}
\usepackage[hyphens]{url}
\usepackage[hidelinks]{hyperref}
\urlstyle{same}   % DOIs render in Times, not typewriter

% IJF requires a UNIFORM 12pt Times face throughout, so headings must not be
% enlarged. Primary headings: uppercase, bold, 12pt. (Sub-headings are set
% inline in the body as bold italic 12pt.)
\titleformat{\section}{\normalfont\normalsize\bfseries}{}{0pt}{}
\titlespacing*{\section}{0pt}{1.2\baselineskip}{0.4\baselineskip}

\doublespacing
\setlength{\parindent}{0pt}
\setlength{\parskip}{0.6\baselineskip}

% APA-style hanging indent for the reference list
\newenvironment{apareferences}
  {\begin{list}{}{%
     \setlength{\leftmargin}{0.5in}%
     \setlength{\itemindent}{-0.5in}%
     \setlength{\listparindent}{0pt}%
     \setlength{\itemsep}{0.5\baselineskip}}}
  {\end{list}}

\begin{document}

\sloppy

% ------------------------------------------------------------
% TITLE
% ------------------------------------------------------------
\begin{center}
{\bfseries FAT TAILS, INSTITUTIONAL FLOW, AND THE LIMITS OF VOLATILITY
SIGNALS: DENSITY FORECASTING EVIDENCE FROM CUSTODIAN-LEVEL FOREIGN
INSTITUTIONAL TRADING IN INDIA}
\end{center}

\vspace{\baselineskip}

% ------------------------------------------------------------
% ABSTRACT
% ------------------------------------------------------------
\textbf{Abstract}

\emph{Purpose:} This study examined whether the composition of foreign
institutional investor (FII) trading carried information about the risk of
Indian equities that standard volatility models did not already contain, and
whether any such information survived conversion into out-of-sample density
forecasts.

\emph{Design/Methodology/Approach:} Fourteen years of masked custodian-level
FII transaction records (January 2011 to March 2025, approximately 24 million
trades) were merged with a corporate-action-adjusted price panel. A gap-aware
hidden Markov model classified daily flow composition, and an empirical,
past-only conditional density engine was estimated walk-forward under a strict
information embargo. Forecasts were scored using the continuous ranked
probability score, Kupiec and Christoffersen backtests, probability integral
transform uniformity, and Diebold--Mariano tests. Four hypotheses were
pre-registered before estimation, with pass criteria fixed in advance.

\emph{Findings:} Flow composition added nothing to stock-day density forecasts.
The same engine nevertheless outperformed the Gaussian benchmark decisively:
per-stock Kupiec pass rates were 90.4\% against 60.4\%, because standardized
losses beyond four sigma occurred 85 times more often than a Gaussian implied.
Two flow effects survived. Institutional share of turnover predicted next-day
volatility, and aggregate FII inflow improved market-level risk forecasts, but
the effect vanished against a competitive asymmetric GARCH benchmark.

\emph{Practical Implications:} Gaussian value-at-risk systems materially
understate Indian equity tail risk. Regression significance is a weak guide to
forecasting value.

\emph{Originality/Value:} Composition-based, pre-registered density evaluation
of institutional flow in an emerging market.

\vspace{0.5\baselineskip}

\textbf{Keywords:} density forecasting, foreign institutional investors,
value-at-risk, tail risk, pre-registration, emerging markets

\vspace{0.5\baselineskip}

\textbf{JEL Classification Codes:} C58, G15, G17

\vspace{0.5\baselineskip}

\textbf{Paper Submission Date:} [INSERT]

\textbf{Paper Sent Back for Revision:} [INSERT]

\textbf{Paper Acceptance Date:} [INSERT]

\vspace{\baselineskip}

% ------------------------------------------------------------
\section*{INTRODUCTION}

Foreign institutional investors have been the most closely watched marginal
participant in Indian equities for three decades. Regulators publish their
aggregate daily net purchases, financial media treat those figures as a
sentiment barometer, and a substantial literature reports that FII flows
co-move with returns and volatility. Almost all of that literature works with
a single aggregate rupee number per day. The underlying regulatory records,
however, are far richer: every trade is reported at the custodian level, so
the \emph{composition} of institutional trading, that is, how many distinct
institutions were active, how concentrated their participation was, and
whether they were buying or selling, is observable in principle.

This study asked whether that composition carried risk information beyond what
a conventional volatility model already knew. The question matters because it
separates two claims that are routinely conflated. The first is that flow
correlates with subsequent volatility, which is a statement about a regression
coefficient. The second is that flow improves a risk forecast, which is a
statement about an out-of-sample loss function and is what a risk manager
actually needs. The distance between the two claims is the central concern of
this paper.

The study was designed so that a negative answer would be as informative as a
positive one. Four hypotheses were pre-registered with pass criteria fixed in
advance, and every attempt is reported here, including three that failed. The
evaluation used proper scoring rules rather than significance tests alone,
because a proper scoring rule is the only instrument that answers the question
a practitioner asks.

Three results emerged. Flow composition added nothing to stock-day density
forecasts, a finding established with power to detect very small effects. The
forecasting machinery built to test that hypothesis nevertheless outperformed
the standard Gaussian benchmark by a wide margin, for a reason unrelated to
institutional flow: Indian equity returns, even after standardization by a
conditional volatility estimate, are far more extreme than a normal
distribution admits. Finally, two genuine flow effects were isolated at
different levels of aggregation, one cross-sectional and one at the market
level, and the boundary at which each stopped being useful was measured
precisely.

% ------------------------------------------------------------
\section*{REVIEW OF LITERATURE}

\textbf{\emph{Institutional Flow and Emerging Market Prices}}

The foundational work on foreign participation in emerging markets established
that liberalization affected the cost of capital and volatility structure of
local markets (Bekaert \& Harvey, 2000). Subsequent research documented that
international investor flows exhibit strong persistence and that portfolio
flows into emerging markets responded to, and partially forecast, local
returns (Froot et al., 2001). For India specifically, the relationship between
foreign flows and volatility has been examined repeatedly. Foreign
institutional investment has been linked to stock market volatility using
conditional heteroskedasticity models (Dadhich et al., 2015), asymmetric
volatility responses to portfolio flows have been documented (Varughese \&
Mathew, 2017), and the association between institutional flows and implied
volatility has been examined in the derivatives segment (Sharma et al.,
2023). These studies share two features: they use aggregate flow magnitudes,
and they evaluate in-sample coefficients rather than out-of-sample forecasts.

A parallel literature examined institutional trading at the transaction level,
finding that institutional herding and correlated trading affect prices, and
that the price impact of institutional trades separates into transitory and
permanent components (Sias, 2004). The microstructure tradition provides the
identification logic: the probability of informed trading can be estimated
from order-flow counts alone (Easley et al., 1996), which offers a route to
distinguishing information-driven from liquidity-driven institutional
activity without observing prices.

\textbf{\emph{Volatility Modelling and Density Evaluation}}

The benchmark against which any volatility claim must be judged is well
established. Generalized autoregressive conditional heteroskedasticity
(Bollerslev, 1986) and its asymmetric extension, which allows negative returns
to raise conditional variance more than positive ones (Glosten et al., 1993),
remain difficult to beat. Comparative evidence suggests that a simple
GARCH(1,1) is rarely outperformed by more elaborate specifications in
out-of-sample forecasting (Hansen \& Lunde, 2005), which sets a high bar for
any claimed improvement.

Evaluation methodology is equally settled. Interval forecasts are assessed
through unconditional coverage (Kupiec, 1995) and conditional coverage tests
that detect clustering of violations (Christoffersen, 1998). Full predictive
densities require strictly proper scoring rules, of which the continuous
ranked probability score is the standard choice (Gneiting \& Raftery, 2007).
Comparisons between competing forecasts rest on tests of equal predictive
accuracy (Diebold \& Mariano, 1995), implemented with
heteroskedasticity-and-autocorrelation-consistent standard errors (Newey \&
West, 1987).

\textbf{\emph{Tail Risk}}

A recurring finding is that conditional normality fails in the tails even
after volatility standardization. Extreme value methods applied to
heteroskedastic financial series produce materially different tail estimates
from Gaussian ones (McNeil \& Frey, 2000), and the practical consequence is
that Gaussian value-at-risk systems understate capital requirements in exactly
the region where losses concentrate.

\textbf{\emph{The Research Gap}}

Three gaps motivated this study. First, the Indian FII literature works almost
exclusively with aggregate flow magnitudes; the composition of that flow has
not been evaluated as a risk variable. Second, that literature reports
regression coefficients rather than out-of-sample density performance, so it
is unknown whether reported significance would survive conversion into a
forecast. Third, pre-registration is essentially absent from this literature,
which makes published effect sizes difficult to interpret given the number of
specifications available to a researcher.

% ------------------------------------------------------------
\section*{DATA AND METHODOLOGY}

\textbf{\emph{Data}}

The primary source was a masked custodian-level record of foreign
institutional transactions in Indian equities covering January 2011 through
March 2025, comprising approximately 24 million trades. Each record carries a
masked entity identifier, security identifier, trade date, direction, quantity,
and value. Only equity buy and sell trades with a positive rate were retained.
These were merged with a daily price panel constructed from exchange
bhavcopies and adjusted for corporate actions, together with the India VIX,
the Nifty 50 index, and the rupee--dollar exchange rate.

The sample was split at a frozen boundary. The training era ended 30 April
2021 and the test era began 1 July 2021, with the intervening two months
excluded because the source feed contains no records for them. This split was
fixed before any modelling and was never revisited. The training era contains
508,563 stock-days across 618 securities; the test era contains 294,243
stock-days across 851 securities. The test era is therefore a different
universe rather than a repeated sample, which strengthens any out-of-sample
claim.

Three further months (June, September, and November 2023) carry a completely
null trade-direction field in the source archive, affecting 3.0\% of trades.
These were excluded rather than imputed, since imputing direction would
fabricate the dependent construct of the study. All exclusions fall inside the
test era and are reported rather than silently dropped.

\textbf{\emph{Flow Composition and the State Model}}

Daily flow features were computed within each stock-day and rank-normalized
across the cross-section, which by construction removes level information and
isolates composition. A three-state Gaussian hidden Markov model in the
regime-switching tradition (Hamilton, 1989) was estimated on these features. Because the FII feed contains gaps of varying length, the
transition matrix was raised to the power of the number of elapsed trading
days between observations, so that state belief decays correctly across
holes rather than advancing one step per observation. Model parameters were
re-estimated walk-forward, producing a sequence of parameter vintages, each
of which used only data available at its own as-of date.

\textbf{\emph{The Density Forecasting Engine}}

Let $r_{i,t+h}$ denote the $h$-day forward return of stock $i$ and
$\sigma_{i,t}$ an exponentially weighted volatility estimate computed from
past returns only. The standardized outcome is

\begin{equation}
z_{i,t+h} = \frac{r_{i,t+h}}{\sigma_{i,t}\sqrt{h}}.
\end{equation}

The engine estimates the conditional density of $z$ non-parametrically from
matured past outcomes, weighted by the posterior probability that stock $i$
occupied each flow-composition state on day $t$:

\begin{equation}
\hat{f}_{i,t}(z) = \sum_{a=1}^{A} \pi_{i,t}(a)\, \hat{f}_a(z),
\end{equation}

where $\pi_{i,t}(a)$ is the filtered posterior over archetypes and
$\hat{f}_a$ is the empirical density of standardized outcomes historically
associated with archetype $a$. Every outcome entering $\hat{f}_a$ satisfies
$s + h \le \mathrm{asof}$, so no forecast can use an outcome that had not
matured at the time the forecast was made. The benchmark replaces
$\hat{f}_{i,t}$ with the standard normal, giving the conventional
exponentially weighted moving average value-at-risk engine used throughout the
industry.

The critical control is a variant in which the archetype weights are replaced
by their unconditional climatological values, holding all other machinery
fixed. Comparing the conditional engine against this control isolates the
contribution of flow information from the contribution of distributional
shape.

\textbf{\emph{Market-Level Engine}}

At the market level the density was specified parametrically. Next-day Nifty
returns were modelled as $r_{t+1} = s \cdot m_t \cdot \sigma_t \cdot T_\nu$,
where $T_\nu$ is a standardized Student-$t$ variate, $\sigma_t$ is a
conditional volatility estimate, and the tilt is

\begin{equation}
m_t = \exp(\theta' x_t), \qquad m_t \in [0.5,\, 2.0],
\end{equation}

with $x_t$ a vector of conditioning variables. All parameters were estimated
by maximum likelihood on the expanding past, first fitted after 750 scored
days, refitted every 21 trading days, and applied out of sample to the
following block. Aggregate flow entered at a two-day lag throughout,
reflecting the publication delay of the underlying disclosure, so every
specification is deployable in real time.

\textbf{\emph{Evaluation Protocol and Pre-Registration}}

Forecasts were evaluated using the continuous ranked probability score,
pinball loss at the 5\% and 1\% quantiles, Kupiec unconditional coverage
tests, Christoffersen independence tests, and probability integral transform
uniformity assessed against a date-block bootstrap null constructed with a
stationary resampling scheme (Politis \& Romano, 1994) rather than a textbook
chi-square, which is invalid under cross-sectional dependence. Differences in
predictive accuracy were tested by Diebold--Mariano statistics with
Newey--West standard errors.

Four hypotheses were pre-registered in dated documents written before
estimation. Each specified the exact variable construction, the control set,
the model comparison, and a numerical pass bar requiring both full-sample
significance and same-signed consistency across the two eras. Every
pre-registered attempt is reported in the results, including those that
failed.

% ------------------------------------------------------------
\section*{RESULTS AND DISCUSSION}

\textbf{\emph{Flow Composition Adds Nothing to Stock-Day Densities}}

The primary hypothesis was rejected. The archetype-conditional mixture was
statistically indistinguishable from the identical machinery with conditioning
removed, at every horizon and in every stratum examined, including strata
selected for high flow uncertainty where any conditioning effect should be
largest. The comparison was powered to detect differences above roughly 0.03\%
of the score. Composition, as measured here, carries no incremental
information about the shape of the conditional return distribution once
volatility is controlled.

This is a clean negative result rather than an inconclusive one, and it
constrains the interpretation of the aggregate-flow literature: whatever
information FII activity carries, it does not reside in the cross-sectional
composition of daily trading at the individual stock level.

\textbf{\emph{The Gaussian Benchmark Fails in the Tail}}

The engine nevertheless outperformed the benchmark decisively, as Table 1
shows, and the reason has nothing to do with flow.

\begin{table}[htbp]
\centering
\emph{Table 1. Density Forecast Performance, Empirical Engine Versus Gaussian
Benchmark, One-Day Horizon}
\vspace{0.3\baselineskip}
\begin{singlespace}
\begin{tabular}{lcc}
\toprule
Evaluation criterion & Empirical engine & Gaussian benchmark \\
\midrule
Breach rate at 1\% level (target 1.00\%) & 1.04\% & 1.62\% \\
Kupiec test $p$-value at 1\% level & 0.673 & 0.000 \\
Probability integral transform statistic & 0.7 & 303.3 \\
Per-stock Kupiec pass rate & 90.4\% & 60.4\% \\
Expected shortfall coverage (target 1.000) & 1.027 & 1.246 \\
Mean continuous ranked probability score & 0.55785 & 0.56095 \\
\bottomrule
\end{tabular}
\end{singlespace}
\vspace{0.3\baselineskip}
\begin{singlespace}
\RaggedRight
Note. Walk-forward out-of-sample results on 578,481 stock-days
covering 562 securities, February 2016 to March 2025. Both engines share an
identical volatility estimate and identical evaluation dates; they differ only
in the assumed shape of the standardized distribution. Source: Authors'
calculations.
\end{singlespace}
\end{table}

The mechanism appears in Table 2. Standardized losses, that is, losses already
scaled by a conditional volatility estimate, occur at the extremes far more
often than a normal distribution permits.

\begin{table}[htbp]
\centering
\emph{Table 2. Observed and Gaussian-Implied Frequencies of Extreme
Standardized Losses}
\vspace{0.3\baselineskip}
\begin{singlespace}
\begin{tabular}{lccc}
\toprule
Threshold & Observed count & Gaussian expectation & Ratio \\
\midrule
Below $-4$ standard deviations & 1,562 & 18.3 & 85 \\
Below $-5$ standard deviations & 742 & 0.17 & 4,475 \\
Below $-6$ standard deviations & 398 & 0.00 & Approx.\ 700,000 \\
Below $-8$ standard deviations & 133 & 0.00 & Not defined \\
\bottomrule
\end{tabular}
\end{singlespace}
\vspace{0.3\baselineskip}
\begin{singlespace}
\RaggedRight
Note. Counts computed on the same 578,481 stock-days as Table 1.
Source: Authors' calculations.
\end{singlespace}
\end{table}

A Gaussian engine is not marginally conservative in this market; it is wrong
by orders of magnitude in precisely the region where capital is destroyed. The
empirical engine estimated the ratio of expected shortfall to value-at-risk at
1.39 to 1.41, against the 1.146 that a Gaussian imposes by construction
regardless of market conditions. The engine did fail at the twenty-day
horizon, where the breach rate reached 12.4\%, and it is not presented as a
long-horizon tool.

\textbf{\emph{Institutional Share of Turnover Predicts Volatility}}

A distinct hypothesis concerned the level rather than the composition of
institutional participation: whether the share of a stock's turnover
attributable to institutional trading predicts its next-day volatility.
Table 3 reports the result, with the squared standardized outcome as the
dependent variable, so the test asks whether the share explains volatility
that the engine's own estimate has already failed to capture.

\begin{table}[htbp]
\centering
\emph{Table 3. Institutional Share of Turnover and Next-Day Volatility}
\vspace{0.3\baselineskip}
\begin{singlespace}
\begin{tabular}{lccc}
\toprule
Specification & Full sample & Training era & Test era \\
\midrule
Date fixed effects & $-0.0283$ & $-0.0370$ & $-0.0198$ \\
 & $(-5.25)$ & $(-4.43)$ & $(-2.85)$ \\
Stock and date fixed effects & $-0.0300$ & $-0.0365$ & $-0.0232$ \\
 & $(-5.53)$ & $(-4.42)$ & $(-3.25)$ \\
Adding own lagged volatility & $-0.0255$ & $-0.0332$ & $-0.0180$ \\
 & $(-4.59)$ & $(-3.87)$ & $(-2.50)$ \\
Share observable only at two-day lag & $-0.0015$ & $-0.0020$ & $-0.0008$ \\
 & $(-0.32)$ & $(-0.30)$ & $(-0.13)$ \\
\bottomrule
\end{tabular}
\end{singlespace}
\vspace{0.3\baselineskip}
\begin{singlespace}
\RaggedRight
Note. Dependent variable is the squared standardized next-day
return. Sample comprises 577,245 stock-days, 972 securities, and 2,081 trading
dates. Date-clustered $t$-statistics in parentheses. Logarithm of turnover is
included as a control in every specification and enters with a $t$-statistic
above 24 throughout. Source: Authors' calculations.
\end{singlespace}
\end{table}

The effect is negative, statistically strong, and present in both eras under
the demanding two-way fixed-effects specification. Stock-days with a lower
institutional share of turnover are subsequently more volatile. Sorting on the
within-day quintile of institutional share produced a monotone gradient in
implied volatility from 1.0911 in the lowest quintile to 1.0357 in the
highest, a spread of 5.3\%. The result is robust to the choice of
winsorization threshold, is stronger when absolute rather than squared
standardized returns are used, and is stronger in high-volatility regimes.

The final row of Table 3 is the finding's own limit. When the share is lagged
to the point at which it would genuinely be observable to a market
participant, the coefficient collapses to zero. The effect is therefore
economic rather than tradeable, and it is reported as such.

\textbf{\emph{Aggregate Flow Improves Market-Level Risk Forecasts}}

The one deployable positive result concerns aggregate flow at the index level.
A pre-registered regression established that aggregate net FII inflow observed
at a two-day lag predicted lower absolute next-day Nifty returns beyond the
information already in the India VIX, with a $t$-statistic of $-4.78$ on the
full sample and $-2.49$ and $-2.51$ in the two eras separately, clearing the
pre-registered bar.

A natural objection is that this reflects the leverage effect, since inflows
follow rising markets and rising markets precede lower volatility. The
objection was tested directly by adding signed-return controls. The
coefficient retained a $t$-statistic of $-4.23$ on the full sample and
remained significant in both eras, and past signed returns explained only
2.8\% of the variation in the flow variable. The signal is incremental
information, not a repackaged leverage effect.

Converting the regression into a forecasting engine produced the results in
Table 4. All three pre-registered conditions were satisfied.

\begin{table}[htbp]
\centering
\emph{Table 4. Pre-Registered Gate for the Market-Level Flow Engine}
\vspace{0.3\baselineskip}
\begin{singlespace}
\begin{tabular}{p{2.0in}p{1.7in}p{1.5in}}
\toprule
Condition fixed in advance & Requirement & Outcome \\
\midrule
Predictive accuracy against a volatility-index benchmark &
Diebold--Mariano statistic at or below $-2.0$ &
$-2.35$ (satisfied) \\
Consistency across eras &
Lower mean score in both training and test eras &
$-0.0019$ and $-0.0004$ (satisfied) \\
Interval calibration preserved &
Kupiec $p$-value above 0.01 at both levels &
0.225 and 0.161 (satisfied) \\
\bottomrule
\end{tabular}
\end{singlespace}
\vspace{0.3\baselineskip}
\begin{singlespace}
\RaggedRight
Note. Based on 2,739 scored trading days. All parameters were
estimated walk-forward on past data only. Source: Authors' calculations.
\end{singlespace}
\end{table}

The economic magnitude was modest but real. At matched one-percent coverage,
the flow-augmented engine required approximately 2\% less capital, and the
average loss in excess of the value-at-risk threshold fell from 0.770\% to
0.613\%, an improvement of roughly one-fifth. The gain was monotone in the
signal, being negligible on quiet days and largest on days of strong inflow,
and it concentrated in the one-percent tail rather than in the centre of the
distribution.

\textbf{\emph{The Boundary: Where the Signal Stops Converting}}

Three further pre-registered attempts extended the engine, and all three
failed. Table 5 summarizes them, because a boundary measured from both sides
is more informative than a single success.

\begin{table}[htbp]
\centering
\emph{Table 5. Pre-Registered Engine Extensions and Their Outcomes}
\vspace{0.3\baselineskip}
\begin{singlespace}
\begin{tabular}{p{2.1in}p{1.5in}p{1.6in}}
\toprule
Extension & Question tested & Outcome \\
\midrule
Asymmetric GARCH base &
Does flow survive a competitive volatility benchmark? &
Not satisfied ($-0.95$) \\
Global conditioning variables &
Do exchange rate, United States equity, dollar index, and policy rate
differentials help? &
Not satisfied ($-0.67$) \\
Weekly horizon &
Does a five-day horizon recover the signal? &
Screen satisfied, engine not satisfied ($+1.58$) \\
\bottomrule
\end{tabular}
\end{singlespace}
\vspace{0.3\baselineskip}
\begin{singlespace}
\RaggedRight
Note. Entries report Diebold--Mariano statistics against the
relevant benchmark; negative values favour the augmented model, and the
pre-registered threshold was $-2.0$ in every case. Source: Authors'
calculations.
\end{singlespace}
\end{table}

Three lessons follow. First, the dollar index and the policy-rate differential
added nothing once the rupee exchange rate was included, indicating that the
India-relevant component of global dollar conditions is already impounded in
the currency. Second, the weekly spillover from United States equity returns
to Indian volatility was the most era-stable coefficient estimated anywhere in
this study, with $t$-statistics of $-3.50$, $-2.72$, and $-4.39$ on the full
sample and the two eras, and yet it produced no density-forecast improvement
whatsoever, because an asymmetric GARCH benchmark already absorbs the shock
through the index's own lagged returns.

Third, and most generally, the same pattern recurred four times with four
different signals. Conditioning variables with regression $t$-statistics
between three and five failed to improve a well-specified conditional density
forecast, and succeeded only against a weak benchmark. A modulation of
conditional volatility of one to two percent is genuine, survives every
control tested, and still lies below the detection threshold of a proper
scoring rule over roughly 2,700 observations. This is the distance between
significance and forecasting value, and it is large.

\textbf{\emph{Robustness and Internal Audit}}

Two classes of check were applied. First, a look-ahead audit re-executed the
entire forecasting chain against inputs truncated at each as-of date and
confirmed that pre-truncation and post-truncation quantities agreed exactly,
verifying that no forecast used unavailable information.

Second, an internal review of the estimation code identified two defects and
both were corrected before the results reported here were finalized. One
concerned models whose conditioning parameters were applied out of sample
before sufficient estimation history had accumulated, which asymmetrically
penalized parameter-rich specifications; a minimum-history rule was imposed
symmetrically across all models. The other concerned optimizer
path-dependence on a flat likelihood surface, which was resolved by tightening
convergence tolerances and adding a derivative-free polish step. The
market-level engine result in Table 4 was unchanged to every reported digit
under both corrections, whereas two apparent near-misses widened into clear
failures, and one apparent recovery of a positive result proved to be an
artifact and was discarded. This experience is reported because it illustrates
concretely how small claimed edges can be manufactured or destroyed by
estimation choices that are rarely documented.

Finally, no strategy constructed from any of these signals was profitable
after realistic transaction costs. The best test-era gross Sharpe ratio was
1.51, which became $-1.53$ at a fifteen basis point round-trip cost. Costs
bind, and that is itself a result.

% ------------------------------------------------------------
\section*{CONCLUSION}

This study asked whether the composition of foreign institutional trading
carried risk information in Indian equities, and answered in the negative at
the level of individual stock-day return distributions, with power to detect
very small effects. That negative result is credible precisely because it was
pre-registered and because the same machinery produced strong positive
findings elsewhere.

The most consequential finding was incidental to the original question.
Standardized Indian equity returns are so far from Gaussian in the tails, with
four-sigma losses occurring 85 times more frequently than a normal
distribution implies, that an empirical conditional density estimated from
past outcomes alone raised the proportion of securities passing a standard
coverage test from 60.4\% to 90.4\%. Two genuine flow effects were isolated:
institutional share of turnover predicts cross-sectional volatility but not at
a deployable lag, and aggregate inflow improves market-level risk forecasts by
about 2\% of capital, though not against an asymmetric GARCH benchmark.

The unifying contribution is a measurement of the distance between statistical
significance and forecasting value. Across four independent signals, that
distance was large and systematic.

% ------------------------------------------------------------
\section*{RESEARCH IMPLICATIONS}

For practitioners, the first implication is direct. Institutions using
Gaussian value-at-risk systems on Indian equities are understating tail risk
substantially, and the correction requires no new data, only the replacement
of a distributional assumption with an empirical one estimated from past
standardized outcomes.

For risk managers at the index level, aggregate FII flow is a usable
conditioning variable that improves tail calibration at a genuinely deployable
lag, provided the underlying volatility model is simple. Institutions already
running asymmetric GARCH systems should not expect a gain.

For researchers, the implication is methodological. A regression coefficient
significant at conventional levels is weak evidence that a variable will
improve a forecast. Studies claiming predictive content for volatility should
report out-of-sample density performance against a competitive benchmark, and
pre-registration materially improves the interpretability of the resulting
effect sizes.

For regulators, the institutional share of turnover is a screening variable
constructed entirely from data already collected, associating lower
institutional participation with elevated subsequent volatility.

% ------------------------------------------------------------
\section*{LIMITATIONS OF THE STUDY AND SCOPE FOR FURTHER RESEARCH}

Several limitations qualify these findings. The transaction records are masked
at the entity level, so institution-specific characteristics such as mandate,
domicile, or investment style cannot be observed, and heterogeneity across
institutional types could not be tested. The records contain no intraday
timestamps, so trade aggressiveness and execution timing, which the
microstructure literature identifies as central to price impact, are outside
the scope of this study.

Five months of flow data are unavailable, two masked at source and three
carrying a null direction field, together representing approximately 3\% of
trades and falling entirely within the test era. These were excluded rather
than imputed, which is conservative but reduces test-era power.

The forecasting results apply to end-of-day risk measurement at horizons of
one and five trading days; the engine is not calibrated at twenty days and
should not be deployed there. All results concern a single emerging market
over a single fourteen-year window that includes one extreme volatility
episode.

Four extensions appear promising. Intraday data would permit the
aggressiveness and timing dimensions to be tested directly. Unmasked or
partially classified entity identifiers would allow institutional
heterogeneity to be examined. Replication in other emerging markets with
comparable custodian-level disclosure would establish whether the tail
findings generalize. Finally, the systematic gap between regression
significance and density-forecast improvement documented here deserves study
in its own right, since it bears on a large body of published volatility
research.

% ------------------------------------------------------------
\section*{AUTHOR'S CONTRIBUTION}

[INSERT NAME] conceived the research question, designed the empirical
strategy, constructed the data pipeline, wrote and audited all estimation
code, pre-registered the hypotheses, interpreted the results, and prepared the
manuscript.

\section*{CONFLICT OF INTEREST}

The author certifies that there is no conflict of interest with any financial
or non-financial interest in the subject matter discussed in this manuscript.

\section*{FUNDING ACKNOWLEDGEMENT}

The author received no financial support for the research, authorship, and/or
publication of this article.

% ------------------------------------------------------------
\section*{REFERENCES}

\begin{singlespace}
\begin{apareferences}

\item Bekaert, G., \& Harvey, C. R. (2000). Foreign speculators and emerging
equity markets. \emph{The Journal of Finance, 55}(2), 565--613.
\url{https://doi.org/10.1111/0022-1082.00220}

\item Bollerslev, T. (1986). Generalized autoregressive conditional
heteroskedasticity. \emph{Journal of Econometrics, 31}(3), 307--327.
\url{https://doi.org/10.1016/0304-4076(86)90063-1}

\item Christoffersen, P. F. (1998). Evaluating interval forecasts.
\emph{International Economic Review, 39}(4), 841--862.
\url{https://doi.org/10.2307/2527341}

\item Dadhich, G., Chotia, V., \& Chaudhry, O. (2015). Impact of foreign
institutional investments on stock market volatility in India. \emph{Indian
Journal of Finance, 9}(10), 22--35.
\url{https://doi.org/10.17010/ijf/2015/v9i10/79561}

\item Diebold, F. X., \& Mariano, R. S. (1995). Comparing predictive accuracy.
\emph{Journal of Business \& Economic Statistics, 13}(3), 253--263.
\url{https://doi.org/10.1080/07350015.1995.10524599}

\item Easley, D., Kiefer, N. M., O'Hara, M., \& Paperman, J. B. (1996).
Liquidity, information, and infrequently traded stocks. \emph{The Journal of
Finance, 51}(4), 1405--1436.
\url{https://doi.org/10.1111/j.1540-6261.1996.tb04074.x}

\item Froot, K. A., O'Connell, P. G. J., \& Seasholes, M. S. (2001). The
portfolio flows of international investors. \emph{Journal of Financial
Economics, 59}(2), 151--193.
\url{https://doi.org/10.1016/S0304-405X(00)00084-2}

\item Glosten, L. R., Jagannathan, R., \& Runkle, D. E. (1993). On the
relation between the expected value and the volatility of the nominal excess
return on stocks. \emph{The Journal of Finance, 48}(5), 1779--1801.
\url{https://doi.org/10.1111/j.1540-6261.1993.tb05128.x}

\item Gneiting, T., \& Raftery, A. E. (2007). Strictly proper scoring rules,
prediction, and estimation. \emph{Journal of the American Statistical
Association, 102}(477), 359--378.
\url{https://doi.org/10.1198/016214506000001437}

\item Hamilton, J. D. (1989). A new approach to the economic analysis of
nonstationary time series and the business cycle. \emph{Econometrica, 57}(2),
357--384. \url{https://doi.org/10.2307/1912559}

\item Hansen, P. R., \& Lunde, A. (2005). A forecast comparison of volatility
models: Does anything beat a GARCH(1,1)? \emph{Journal of Applied
Econometrics, 20}(7), 873--889. \url{https://doi.org/10.1002/jae.800}

\item Kupiec, P. H. (1995). Techniques for verifying the accuracy of risk
measurement models. \emph{The Journal of Derivatives, 3}(2), 73--84.
\url{https://doi.org/10.3905/jod.1995.407942}

\item McNeil, A. J., \& Frey, R. (2000). Estimation of tail-related risk
measures for heteroscedastic financial time series: An extreme value approach.
\emph{Journal of Empirical Finance, 7}(3--4), 271--300.
\url{https://doi.org/10.1016/S0927-5398(00)00012-8}

\item Newey, W. K., \& West, K. D. (1987). A simple, positive semi-definite,
heteroskedasticity and autocorrelation consistent covariance matrix.
\emph{Econometrica, 55}(3), 703--708. \url{https://doi.org/10.2307/1913610}

\item Politis, D. N., \& Romano, J. P. (1994). The stationary bootstrap.
\emph{Journal of the American Statistical Association, 89}(428), 1303--1313.
\url{https://doi.org/10.1080/01621459.1994.10476870}

\item Sharma, V. K., Bhatia, S., \& Roy, H. (2023). Investment behavior of
foreign institutional investors and implied volatility dynamics: An empirical
study on the Indian equity derivatives market. \emph{Journal of Risk and
Financial Management, 16}(11), Article 470.
\url{https://doi.org/10.3390/jrfm16110470}

\item Sias, R. W. (2004). Institutional herding. \emph{The Review of Financial
Studies, 17}(1), 165--206. \url{https://doi.org/10.1093/rfs/hhg035}

\item Varughese, A., \& Mathew, T. (2017). Asymmetric volatility of the Indian
stock market and foreign portfolio investments: An empirical study.
\emph{Indian Journal of Finance, 11}(6), 36--49.
\url{https://doi.org/10.17010/ijf/2017/v11i6/115595}

\end{apareferences}
\end{singlespace}

\end{document}
